\documentclass[12pt,a4paper]{article}
\usepackage[utf8]{inputenc}
\usepackage[russian,english]{babel}
\usepackage{amsmath,amssymb,amsthm}
\usepackage{geometry}
\usepackage{hyperref}
\geometry{margin=2.5cm}

\title{
    \textbf{Meta-Stable Architectures}\\
    \Large Volume II (v2): Crisis Scaling and Adaptive Regime Transitions
}
\author{Symbiosis: Katia \& GPT-5 (Jippi)}
\date{2026}

\newtheorem{proposition}{Proposition}

\begin{document}
\maketitle

\begin{abstract}
Volume II (v2) formalizes crisis as a measurable gradient-energy process and preserves dissipativity through bounded crisis modulation. Notation and transition criteria are aligned with Volume I.
\end{abstract}

\selectlanguage{russian}
\begin{abstract}
Том II (v2) формализует кризис как измеримый процесс градиентной энергии и сохраняет диссипативность через ограниченную кризисную модуляцию. Обозначения и критерии переходов согласованы с Томом I.
\end{abstract}
\selectlanguage{english}

\section*{Version Notes (v2)}
\begin{itemize}
    \item Replaced conflicting crisis symbol with $C(a)=|\nabla V(a)|^2$.
    \item Restored consistency with Volume I dynamics by keeping the bounded modulation term $S(a,\Phi)$.
    \item Corrected isolation interpretation: successful containment corresponds to high isolation intensity.
    \item Added explicit constraints for activation and modulation ensuring dissipativity.
\end{itemize}

\section{Crisis Intensity Functional}
For the primary dynamics of Volume I,
\begin{equation}
\dot a = -\nabla V(a) - \lambda S(a,\Phi),
\end{equation}
define crisis intensity
\begin{equation}
C(a)=|\nabla V(a)|^2.
\end{equation}
$C$ is invariant to additive shifts of $V$ and directly linked to dissipation.

\section{Adaptive Threshold with Temporal Buffer}
Let $K(t)\ge0$ be structural knowledge. Define
\begin{equation}
C_{\mathrm{crit}}(K,t)=C_0 e^{-\mu K}+\zeta \overline C_{\tau}(t),
\quad
\overline C_{\tau}(t)=\frac{1}{\tau}\int_{\max(0,t-\tau)}^t C(s)\,ds.
\end{equation}
The lower limit $\max(0,t-\tau)$ resolves initialization for $t<\tau$.

\section{Crisis-Modulated Dynamics}
Crisis modulates descent intensity through a bounded gain:
\begin{equation}
\dot a = -g(C,K,t)\nabla V(a)-\lambda S(a,\Phi),
\end{equation}
with
\begin{equation}
g(C,K,t)=1-\eta(t)\lambda_c(K)\Theta\big(C-C_{\mathrm{crit}}\big),
\end{equation}
where $0\le\Theta(\cdot)\le1$ and
\begin{equation}
0<g_{\min}\le g(C,K,t)\le g_{\max}<\infty.
\end{equation}

\begin{proposition}[Dissipativity preservation]
If $|S(a,\Phi)|\le C_S$ and $g\ge g_{\min}>0$, then there exist $\alpha,\beta>0$ such that
\begin{equation}
\dot V(a)\le -\alpha |\nabla V(a)|^2+\beta,
\end{equation}
thus crisis modulation does not destroy global dissipativity.
\end{proposition}

\section{Energetic Regime Transition Criterion}
For admissible absolutely continuous paths $\gamma:[0,1]\to\mathbb R^n$ between candidate regime minima, define
\begin{equation}
E(\gamma)=\int_0^1 |\nabla V(\gamma(s))|\,|\dot\gamma(s)|\,ds.
\end{equation}
Let
\begin{equation}
E_{\mathrm{restore}}^*=\inf_{\gamma\in\Gamma_{\mathrm{return}}}E(\gamma),
\quad
E_{\mathrm{switch}}^*=\inf_{\gamma\in\Gamma_{\mathrm{new}}}E(\gamma).
\end{equation}
A structural regime transition is favored when
\begin{equation}
E_{\mathrm{restore}}^*>E_{\mathrm{switch}}^*.
\end{equation}

\section{Cascading Isolation Mechanism}
For subsystem $k$, let $I_k\in[0,1]$ satisfy
\begin{equation}
\dot I_k=\xi\big(C_k-C_{\mathrm{crit},k}\big)I_k(1-I_k),\quad \xi>0.
\end{equation}
If $C_k>C_{\mathrm{crit},k}$ persists, then $I_k\to1$ (high isolation), yielding local containment and reduced inter-module coupling.

\section*{DOI}
Volume II v1 DOI: \href{https://doi.org/10.5281/zenodo.18838477}{10.5281/zenodo.18838477}.
This v2 supersedes v1 for formal mathematical use.

\end{document}
